\documentclass[aspectratio=169]{beamer} % Formato widescreen
\usetheme{Madrid} % Tema limpo e profissional
\usecolortheme{default}

% Pacotes para português e gráficos
\usepackage[utf8]{inputenc}
\usepackage[T1]{fontenc}
\usepackage[portuguese]{babel}
\usepackage{graphicx}
\usepackage{booktabs}

% Informações da Capa
\title[Projeto EcoTátil]{EcoTátil: Alfabetização em Braille Multimodal}
\subtitle{Pitch para Investidores-Anjos - Fase 3}
\author[Jumpers]{Marcelo Arbage, Theo Mello, Thiago Araujo e Rafael Collares}
\institute[UFPA]{
  Empresa Jumpers - Tecnologia Assistiva\\
  Universidade Federal do Pará (UFPA)
}
\date{17 de Julho de 2026}

\begin{document}

% Slide 1: Capa
\begin{frame}
  \titlepage
\end{frame}

% Slide 2: Índice
\begin{frame}{Agenda}
  \tableofcontents
\end{frame}

% Seção 1: O Problema
\section{O Problema}
\begin{frame}{O Desafio da Alfabetização Inclusiva}
  \begin{itemize}
    \item \textbf{A Dor:} Pessoas com deficiência visual enfrentam grandes barreiras na alfabetização tátil.
    \item \textbf{O Gap de Mercado:} O uso exclusivo de leitores de tela não consolida habilidades de leitura e escrita, afetando o desenvolvimento cognitivo.
    \item \textbf{Custo:} Dispositivos de interface Braille (displays) costumam ser extremamente caros e inacessíveis para o público geral e escolas públicas.
  \end{itemize}
\end{frame}

% Seção 2: A Solução
\section{A Solução: EcoTátil}
\begin{frame}{Conheça o EcoTátil}
  \begin{itemize}
    \item \textbf{O que é:} Uma célula Braille eletrônica interativa de baixo custo.
    \item \textbf{Diferencial - Aprendizado Multimodal:} Combina resposta tátil imediata com feedback auditivo.
    \item \textbf{Como funciona:} O usuário insere um caractere, o sistema processa a informação, eleva os pinos formando a letra em Braille e emite o som correspondente.
  \end{itemize}
  % \begin{figure}
  %   \includegraphics[width=0.4\textwidth]{sua_imagem_do_gabinete_3D.png}
  %   \caption{Gabinete impresso em 3D}
  % \end{figure}
\end{frame}

% Seção 3: Viabilidade Técnica
\section{Viabilidade Técnica (MVP)}
\begin{frame}{Arquitetura e Componentes}
  \begin{columns}
    \column{0.5\textwidth}
    \textbf{Hardware Estratégico}
    \begin{itemize}
      \item Microcontrolador \textbf{ESP32-S} (substituindo o Arduino por conectividade nativa).
      \item 6 Micro servomotores independentes com placa PCA9685.
      \item Módulo DFPlayer Mini para o feedback sonoro.
    \end{itemize}
    
    \column{0.5\textwidth}
    \textbf{Validação}
    \begin{itemize}
      \item Simulação de circuito e firmware validada na plataforma \textbf{Wokwi}.
      \item Placa de Circuito Impresso (PCI) própria projetada no \textbf{KiCad}.
      \item Design ergonômico impresso em 3D.
    \end{itemize}
  \end{columns}
\end{frame}

% Seção 4: Custos e Investimento
\section{Custos e Modelo de Negócio}
\begin{frame}{Estimativa de Custos do Protótipo}
  \begin{itemize}
    \item O projeto foi concebido com foco na viabilidade econômica usando componentes de prateleira (off-the-shelf).
  \end{itemize}
  
  \vspace{0.5cm}
  \begin{table}[]
      \centering
      \begin{tabular}{lc}
           \toprule
           \textbf{Categoria} & \textbf{Custo Estimado (R\$)} \\
           \midrule
           Processamento (ESP32-S) & 55,00 \\
           Atuação (Servos + Driver) & 120,00 \\
           Áudio (DFPlayer + Falante) & 30,00 \\
           Alimentação e Estrutura & 202,00 \\
           \midrule
           \textbf{Custo Total de Produção (Unidade)} & \textbf{407,00} \\
           \bottomrule
      \end{tabular}
  \end{table}
\end{frame}

% Seção 5: Futuro do Produto
\section{Roadmap e Visão de Futuro}
\begin{frame}{Por que investir na Jumpers?}
  \begin{itemize}
    \item \textbf{Evolução Escalável:} Transição de uma célula única para múltiplas células, formando palavras completas.
    \item \textbf{Reconhecimento de Voz Offline:} Substituição para o ESP32-S3, rodando Inteligência Artificial embarcada (Edge AI) sem depender de nuvem.
    \item \textbf{Aprimoramento Mecânico:} Substituição de servos por atuadores piezoelétricos, diminuindo tamanho e ruído.
    \item \textbf{Impacto Social e Lucratividade:} Um mercado inexplorado que une acessibilidade e tecnologia escalável.
  \end{itemize}
\end{frame}

% Slide de Encerramento
\begin{frame}
  \centering
  \Huge \textbf{Obrigado!}\\
  \vspace{1cm}
  \Large \textbf{Equipe Jumpers}\\
  \normalsize jumpersltda@gmail.com
\end{frame}

\end{document}
